% GENERATED FILE. Edit canonical lessons, then run npm run generate:books.
% canonical-source-hash: fnv1a64:b1825ab972583843
% canonical-lessons: ZH-C211-bendi, ZH-C211-waidi, ZH-C211-shizhongxin, ZH-C211-ditu, ZH-C211-daochu, ZH-R211-review-206-211

\chapter{Local or From Out of Town}
\label{ch:zh-local-or-from-out-of-town}
\hlchaptermodality{211}

\begin{chapteropening}
\textbf{By the end of this chapter:} \emph{I can say who is local and who comes from elsewhere, and talk about the city centre and a map.}

Complete the last lesson of chapter 211: I can say who is local and who comes from elsewhere, and talk about the city centre and a map.
\end{chapteropening}

\section[\texorpdfstring{\zh{本地} (běndì)}{本地 (běndì)}]{\zh{本地} (běndì) — local}

\label{lesson:ZH-C211-bendi}



\begin{quote}
Before the new one: say the Mandarin for fitting room, then the Mandarin for to choose.
\end{quote}

\subsection*{You'll want to know: \zh{本地}}
\textbf{\zh{本地}} — \emph{běndì} — \textquotedblleft{}local\textquotedblright{}.

\begin{quote}
\textbf{\zh{我们找一个本地人问路}} \emph{wǒmen zhǎo yí ge běndì rén wènlù} — Let's find someone local and ask the way.
\end{quote}

\begin{cousinweb}
\textbf{\zh{本地}} (\emph{běndì}) is built from \zh{本} 'root, one's own' and \zh{地} 'place', both inherited from Middle Chinese, \emph{pwonX} and \emph{dijH} (Guangyun, in Baxter's transcription), which glosses the first as the root of a tree and 'the beginning'.
\end{cousinweb}

\begin{grammarlens}[title={What you've built}]
One more word for asking where things are.
\end{grammarlens}

\subsection*{Your turn}
\begin{itemize}
\raggedright
  \item \emph{Say it:} \emph{běndì}
  \item \emph{Say it:} \emph{běndì}, once more
  \item \emph{Say it:} \emph{wǒmen zhǎo yí ge běndì rén wènlù} — Let's find someone local and ask the way.
  \item \emph{Recall:} say the Mandarin for fitting room, then the Mandarin for to choose, then say \emph{běndì} again
\end{itemize}

\begin{tcolorbox}[breakable,colback=teal!4,colframe=teal!35!black,title={Before you move on}]
What does \zh{本地} mean? (\textquotedblleft{}Local\textquotedblright{}.) Say it once more.
\end{tcolorbox}
\section[\texorpdfstring{\zh{外地} (wàidì)}{外地 (wàidì)}]{\zh{外地} (wàidì) — from out of town; another part of the country}

\label{lesson:ZH-C211-waidi}



\begin{quote}
Before the new one: say the Mandarin for to choose, then the Mandarin for local.
\end{quote}

\subsection*{You'll want to know: \zh{外地}}
\textbf{\zh{外地}} — \emph{wàidì} — \textquotedblleft{}from out of town; another part of the country\textquotedblright{}.

\begin{quote}
\textbf{\zh{他是外地人每天开车一百公里}} \emph{tā shì wàidì rén měitiān kāichē yìbǎi gōnglǐ} — He's from out of town and drives a hundred kilometres every day.
\end{quote}

\begin{cousinweb}
\textbf{\zh{外地}} (\emph{wàidì}) is built from \zh{外} 'outside' and \zh{地} 'place', both inherited from Middle Chinese, \emph{ngwajH} and \emph{dijH} (Guangyun, in Baxter's transcription).
\end{cousinweb}

\begin{grammarlens}[title={What you've built}]
One more word for asking where things are.
\end{grammarlens}

\subsection*{Your turn}
\begin{itemize}
\raggedright
  \item \emph{Say it:} \emph{wàidì}
  \item \emph{Say it:} \emph{wàidì}, once more
  \item \emph{Say it:} \emph{tā shì wàidì rén měitiān kāichē yìbǎi gōnglǐ} — He's from out of town and drives a hundred kilometres every day.
  \item \emph{Recall:} say the Mandarin for to choose, then the Mandarin for local, then say \emph{wàidì} again
\end{itemize}

\begin{tcolorbox}[breakable,colback=teal!4,colframe=teal!35!black,title={Before you move on}]
What does \zh{外地} mean? (\textquotedblleft{}From out of town; another part of the country\textquotedblright{}.) Say it once more.
\end{tcolorbox}
\section[shìzhōngxīn]{shìzhōngxīn — city centre}

\label{lesson:ZH-C211-shizhongxin}



\begin{quote}
Before the new one: say the Mandarin for local, then the Mandarin for from out of town.
\end{quote}

\subsection*{You'll want to know: shìzhōngxīn}
\emph{shìzhōngxīn} — \textquotedblleft{}city centre\textquotedblright{}.

\begin{quote}
\emph{wǒ jiā zài shìzhōngxīn} — My home is in the city centre.
\end{quote}

Its characters come later; for now it is a word for the ear.

\begin{cousinweb}
\emph{shìzhōngxīn} is built from \emph{shì} 'a market, a city', \zh{中} 'middle' and \zh{心} 'heart', all inherited from Middle Chinese, \emph{dzyiX}, \emph{trjuwng} and \emph{sim} (Guangyun, in Baxter's transcription), which glosses the first as the place where buying and selling is done.
\end{cousinweb}

\begin{grammarlens}[title={What you've built}]
One more word for asking where things are.
\end{grammarlens}

\subsection*{Your turn}
\begin{itemize}
\raggedright
  \item \emph{Say it:} \emph{shìzhōngxīn}
  \item \emph{Say it:} \emph{shìzhōngxīn}, once more
  \item \emph{Say it:} \emph{wǒ jiā zài shìzhōngxīn} — My home is in the city centre.
  \item \emph{Recall:} say the Mandarin for local, then the Mandarin for from out of town, then say \emph{shìzhōngxīn} again
\end{itemize}

\begin{tcolorbox}[breakable,colback=teal!4,colframe=teal!35!black,title={Before you move on}]
What does shìzhōngxīn mean? (\textquotedblleft{}City centre\textquotedblright{}.) Say it once more.
\end{tcolorbox}
\section[dìtú]{dìtú — map}

\label{lesson:ZH-C211-ditu}



\begin{quote}
Before the new one: say the Mandarin for from out of town, then the Mandarin for city centre.
\end{quote}

\subsection*{You'll want to know: dìtú}
\emph{dìtú} — \textquotedblleft{}map\textquotedblright{}.

\begin{quote}
\emph{nǐ yǒu Běijīng de dìtú ma} — Have you got a map of Beijing?
\end{quote}

Its characters come later; for now it is a word for the ear.

\begin{cousinweb}
\emph{dìtú} is built from \zh{地} 'ground' and \emph{tú} 'a plan, a chart', both inherited from Middle Chinese, \emph{dijH} and \emph{du} (Guangyun, in Baxter's transcription), which glosses the second 'to plan'.
\end{cousinweb}

\begin{grammarlens}[title={What you've built}]
One more word for asking where things are.
\end{grammarlens}

\subsection*{Your turn}
\begin{itemize}
\raggedright
  \item \emph{Say it:} \emph{dìtú}
  \item \emph{Say it:} \emph{dìtú}, once more
  \item \emph{Say it:} \emph{nǐ yǒu Běijīng de dìtú ma} — Have you got a map of Beijing?
  \item \emph{Recall:} say the Mandarin for from out of town, then the Mandarin for city centre, then say \emph{dìtú} again
\end{itemize}

\begin{tcolorbox}[breakable,colback=teal!4,colframe=teal!35!black,title={Before you move on}]
What does dìtú mean? (\textquotedblleft{}Map\textquotedblright{}.) Say it once more.
\end{tcolorbox}
\section[dàochù]{dàochù — everywhere}

\label{lesson:ZH-C211-daochu}



\begin{quote}
Before the new one: say the Mandarin for city centre, then the Mandarin for map.
\end{quote}

\subsection*{You'll want to know: dàochù}
\emph{dàochù} — \textquotedblleft{}everywhere\textquotedblright{}.

\begin{quote}
\emph{shìzhōngxīn dàochù dōu shì rén} — There are people everywhere in the city centre.
\end{quote}

Its characters come later; for now it is a word for the ear.

\begin{cousinweb}
\emph{dàochù} is built from \zh{到} 'to reach' and \emph{chù} 'a place', both inherited from Middle Chinese, \emph{tawH} and \emph{tsyhoH} (Guangyun, in Baxter's transcription), which glosses the second 'a place'.
\end{cousinweb}

\begin{grammarlens}[title={What you've built}]
One more word for asking where things are.
\end{grammarlens}

\subsection*{Your turn}
\begin{itemize}
\raggedright
  \item \emph{Say it:} \emph{dàochù}
  \item \emph{Say it:} \emph{dàochù}, once more
  \item \emph{Say it:} \emph{shìzhōngxīn dàochù dōu shì rén} — There are people everywhere in the city centre.
  \item \emph{Recall:} say the Mandarin for city centre, then the Mandarin for map, then say \emph{dàochù} again
\end{itemize}

\begin{tcolorbox}[breakable,colback=teal!4,colframe=teal!35!black,title={Before you move on}]
What does dàochù mean? (\textquotedblleft{}Everywhere\textquotedblright{}.) Say it once more.
\end{tcolorbox}
\section[(dialogue)]{Looking back: from Upstairs to Local or From Out of Town}

\label{lesson:ZH-R211-review-206-211}



\begin{quote}
Say the Mandarin for a building of several floors and storey.
\end{quote}

\subsection*{Your turn}
\begin{itemize}
\raggedright
  \item \emph{Say it:} a building of several floors, storey, to go upstairs, office, reception desk
  \item \emph{Say it:} venue, a plan, to arrange, to take a photo, tourist sight
  \item \emph{Say it:} skilled tradesman, to examine, to install, tool, battery
  \item \emph{Say it:} regulation, safe, to keep up, to disturb, to permit
  \item \emph{Say it:} purple, brown, style of clothes, fitting room, to choose
  \item \emph{Say it:} local, from out of town, city centre, map, everywhere
  \item \emph{Read it:} \zh{地} — each one aloud, then a word that holds it
\end{itemize}

\begin{tcolorbox}[breakable,colback=teal!4,colframe=teal!35!black,title={Before you move on}]
A building of several floors? (\emph{lóu}.) Storey? (\emph{céng}.) To go upstairs? (\emph{shànglóu}.) Office? (\emph{bàngōngshì}.) Reception desk? (\emph{qiántái}.) Venue? (\textbf{\zh{地点}}, \emph{dìdiǎn}.) A plan? (\emph{jìhuà}.) To arrange? (\emph{ānpái}.) To take a photo? (\emph{zhàoxiàng}.) Tourist sight? (\emph{jǐngdiǎn}.) Skilled tradesman? (\emph{shīfu}.) To examine? (\emph{jiǎnchá}.) To install? (\emph{ānzhuāng}.) Tool? (\emph{gōngjù}.) Battery? (\emph{diànchí}.) Regulation? (\emph{guīdìng}.) Safe? (\emph{ānquán}.) To keep up? (\emph{bǎochí}.) To disturb? (\emph{dǎrǎo}.) To permit? (\emph{yǔnxǔ}.) Purple? (\emph{zǐsè}.) Brown? (\emph{zōngsè}.) Style of clothes? (\emph{kuǎnshì}.) Fitting room? (\textbf{\zh{试衣间}}, \emph{shìyījiān}.) To choose? (\emph{xuǎn}.) Local? (\textbf{\zh{本地}}, \emph{běndì}.) From out of town? (\textbf{\zh{外地}}, \emph{wàidì}.) City centre? (\emph{shìzhōngxīn}.) Map? (\emph{dìtú}.) Everywhere? (\emph{dàochù}.)
\end{tcolorbox}
